% Standalone XeLaTeX source. Original: 2026ap_exam.pdf, Japanese PDF pp. 2--3.
\documentclass[11pt,a4paper]{article}
\usepackage[margin=24mm]{geometry}
\usepackage{fontspec}
\usepackage{xeCJK}
\setCJKmainfont{HaranoAjiMincho-Regular.otf}
\setCJKsansfont{HaranoAjiGothic-Medium.otf}
\usepackage{amsmath,amssymb}
\setlength{\parindent}{1em}
\setlength{\parskip}{0.35em}
\title{物理工学専攻 院試解答集 ver2\\2026年 物理学 第1問}
\author{政岡凜太郎}
\date{}
\begin{document}
\maketitle

\section*{図1の幾何と運動の向き}
円錐の母線に沿って頂点から測った距離が $r$ である。図1の上向き $z$ 軸を使うと、A の水平半径と高さ、B の高さはそれぞれ
\begin{equation}
 R_A=r\sin\theta_0,\qquad z_A=r\cos\theta_0,\qquad
 z_B=-(l-r)=r-l. \label{eq:geometry}
\end{equation}
糸が張っているので、A が斜面を外へ進むと B も上昇する。以下では図の円錐に合わせて $0<\theta_0<\pi/2$、また $0<r<l$ の運動を考える。初速度は方位角方向だけなので、$\dot r(0)=0$、$r_0\sin\theta_0\,\dot\varphi(0)=v_0$ である。

\section*{[1] ラグランジアン}
A の速さの二乗は $\dot r^{2}+(r\sin\theta_0)^2\dot\varphi^{2}$、B の速さの二乗は $\dot z_B^{2}=\dot r^{2}$ である。位置エネルギーの基準を頂点の高さに取ると
\[
 T_A=\frac m2\bigl(\dot r^{2}+r^2\sin^2\theta_0\,\dot\varphi^{2}\bigr),\quad
 T_B=\frac M2\dot r^{2},\quad
 U_A=mg r\cos\theta_0,\quad U_B=Mg(r-l).
\]
従って、基準点による定数項も含めた答えは
\begin{equation}
 \boxed{L=\frac{m+M}{2}\dot r^{2}
 +\frac m2r^2\sin^2\theta_0\,\dot\varphi^{2}
 -g(m\cos\theta_0+M)r+Mgl.} \label{eq:L}
\end{equation}

\section*{[2] $r$ の運動方程式}
式\eqref{eq:L}から $\partial L/\partial\dot r=(m+M)\dot r$、$\partial L/\partial r=mr\sin^2\theta_0\,\dot\varphi^{2}-g(m\cos\theta_0+M)$ である。従って
\begin{equation}
 \boxed{(m+M)\ddot r=mr\sin^2\theta_0\,\dot\varphi^{2}
 -g(m\cos\theta_0+M).} \label{eq:radial}
\end{equation}
右辺の第一項は方位角運動が外向きに寄与する項、第二項は二つの質点を同時に上げることに伴う重力項である。

\section*{[3] $z$ 軸まわりの角運動量}
$L$ は $\varphi$ を含まないので、$\frac{d}{dt}(\partial L/\partial\dot\varphi)=0$。A の水平半径 $R_A$ と水平接線方向の速さ $R_A\dot\varphi$ の積から、保存量は
\begin{equation}
 L_z=mr^2\sin^2\theta_0\,\dot\varphi
 =mr_0\sin\theta_0\,v_0. \label{eq:Lz}
\end{equation}
B は $z$ 軸上で動くため、この軸まわりの角運動量を持たない。

\section*{[4] 方位角を消去した方程式}
式\eqref{eq:Lz}より $\dot\varphi=v_0r_0/(r^2\sin\theta_0)$。これを式\eqref{eq:radial}に代入すると
\begin{equation}
 \boxed{(m+M)\ddot r=\frac{mr_0^2v_0^2}{r^3}
 -g(m\cos\theta_0+M).} \label{eq:reduced}
\end{equation}

\section*{[5] 径方向の速さ}
式\eqref{eq:reduced}に $\dot r$ を掛ける。$r^{-3}\dot r=-\frac12\frac{d}{dt}(r^{-2})$ であり、初期条件 $r=r_0,\dot r=0$ を使って積分すると
\begin{align}
 \frac{m+M}{2}\dot r^{2}
 &=\frac{m v_0^2}{2}\left(1-\frac{r_0^2}{r^2}\right)
 -g(m\cos\theta_0+M)(r-r_0),\\
 \dot r^{2}
 &=\boxed{\frac{m v_0^2}{m+M}\left(1-\frac{r_0^2}{r^2}\right)
 -\frac{2g(m\cos\theta_0+M)}{m+M}(r-r_0).} \label{eq:rdotsq}
\end{align}
これは $r=r_0$ で確かに零になる。$l$ は糸長の制約を定めるが、$r$ に依存する運動エネルギー・力には現れない。

\section*{[6] 他方の転回点}
$r_1\ne r_0$ は式\eqref{eq:rdotsq}のもう一つの正の零点である。$r=r_1$ とおき、$1-r_0^2/r_1^2=(r_1-r_0)(r_1+r_0)/r_1^2$ を用いて $r_1-r_0$ を除くと
\[
 2g(m\cos\theta_0+M)r_1^2-mv_0^2r_1-mv_0^2r_0=0.
\]
左辺の定数項は負なので二根の符号は異なり、物理的な根は正の根である。よって
\begin{equation}
 \boxed{r_1=\frac{mv_0^2+\sqrt{m^2v_0^4+8m v_0^2r_0g(m\cos\theta_0+M)}}{4g(m\cos\theta_0+M)}.} \label{eq:r1}
\end{equation}
$r_1>r_0$ か $r_1<r_0$ かは初期加速度の符号、すなわち $mv_0^2/r_0-g(m\cos\theta_0+M)$ の符号で決まる。等号の場合は $r_1=r_0$ の定常回転となり、この設問の異なる二つの転回点にはならない。問題文が想定する往復運動では、$r_0,r_1$ とその間の $r$ は $0<r<l$ を満たす。

\section*{[7] 定常回転を与える半頂角}
ここからは問題文の通り $\theta_0$ だけを変え、$r=r_0$ を保つ別の運動を考える。式\eqref{eq:reduced}で $\ddot r=0$ とすると $mv_0^2/r_0=g(m\cos\theta_0+M)$。従って
\begin{equation}
 \boxed{\theta_0=\arccos\!\left(\frac{v_0^2}{gr_0}-\frac{M}{m}\right).} \label{eq:theta}
\end{equation}
図の非退化な上向き円錐については $0<v_0^2/(gr_0)-M/m<1$ が必要である。

\section*{[8] 微小変位後の角速度}
瞬間的な力は $r$ 方向だけなので $z$ 軸まわりのトルクはなく、式\eqref{eq:Lz}の値は前後で同じである。$r=r_0+\rho$ を代入して
\begin{equation}
 \boxed{\dot\varphi=\frac{v_0r_0}{(r_0+\rho)^2\sin\theta_0}
 =\frac{v_0}{r_0\sin\theta_0}\left(1+\frac{\rho}{r_0}\right)^{-2}.} \label{eq:phirho}
\end{equation}
ここでは後の運動の $\dot\rho$ の初期値は指定しない。

\section*{[9] $\rho$ の一次の運動方程式}
式\eqref{eq:reduced}に $r=r_0+\rho$ を入れ、$|\rho|\ll r_0$ で $(1+\rho/r_0)^{-3}\simeq1-3\rho/r_0$ と展開する。定数項 $mv_0^2/r_0-g(m\cos\theta_0+M)$ は式\eqref{eq:theta}の定常条件で零なので
\begin{equation}
 \boxed{\ddot\rho+\omega^2\rho=0,\qquad
 \omega^2=\frac{3m v_0^2}{(m+M)r_0^2}
 =\frac{3g(m\cos\theta_0+M)}{(m+M)r_0}.} \label{eq:rho}
\end{equation}
$\omega^2>0$ なので、微小なずれに対する復元運動である。求めた式は $\rho/r_0$ の一次までで、$\rho^2$ 以上を捨てている。

\section*{[10] B の運動}
糸の長さから $z_B=r-l=r_0-l+\rho$ である。従って B は平均位置 $z_B=r_0-l$ のまわりを、A の径方向のずれと同じ $\rho$ だけ上下する。式\eqref{eq:rho}より、その運動は鉛直方向の単振動であり、角振動数は $\omega$、周期は
\[
 \boxed{T=\frac{2\pi}{\omega}
 =2\pi\sqrt{\frac{(m+M)r_0^2}{3m v_0^2}}.}
\]
微小力の力積は指定されていないので、振幅と位相はこの設問だけでは定まらない。A が外側へ動くとき B は上昇する。

\end{document}
